\clearpage
\section{Podsumowanie}
\label{sec:podsumowanie}

Celem pracy było sprawdzenie, czy agent wykorzystujący algorytm Deep Q-Network może skutecznie sterować sygnalizacją świetlną na pojedynczym skrzyżowaniu i osiągać lepsze rezultaty niż wybrane klasyczne metody sterowania. Badanie przeprowadzono w środowisku symulacyjnym na modelu skrzyżowania ulicy Belgradzkiej i alei Komisji Edukacji Narodowej. Zakres pracy dobrano tak, aby możliwe było przeprowadzenie kompletnego, kontrolowanego eksperymentu, obejmującego przygotowanie środowiska, trening, walidację oraz ocenę na odrębnym zbiorze testowym.

\subsection{Zrealizowane prace}
\label{subsec:podsumowanie-prace}

W ramach pracy przygotowano model skrzyżowania w symulatorze SUMO, obejmujący geometrię dróg, trzynaście pasów wlotowych, relacje skrętne oraz trzy fazy sygnalizacji. Zaimplementowano dwa punkty odniesienia: sterowanie stałoczasowe oraz sterowanie akomodacyjne typu \textit{gap-based}. Parametry obu metod wyznaczono na podstawie danych treningowych, a następnie zamrożono przed oceną modelu DQN.

Opracowano rozdzielne zbiory treningowy, walidacyjny i testowy, reprezentujące lekkie, średnie, wysokie oraz zmienne natężenie ruchu. Przygotowano dwa warianty danych treningowych, różniące się regularnością napływu pojazdów. Zbiór testowy obejmował dwadzieścia scenariuszy niewykorzystywanych podczas uczenia ani doboru konfiguracji.

Na podstawie biblioteki \texttt{sumo-rl} zbudowano własną warstwę integrującą SUMO z implementacją algorytmu DQN dostępną w bibliotece Stable-Baselines3. Dostosowano przebieg kroku środowiska, obsługę wielu plików tras oraz sposób zbierania metryk, dzięki czemu wszystkie metody oceniano według wspólnej procedury. Rejestrowano średnie opóźnienie, średni czas oczekiwania, długość kolejek, przepustowość i częstotliwość przełączeń sygnalizacji.

Proces uczenia obejmował ocenę kolejnych checkpointów, porównanie obu wariantów zbioru treningowego oraz sprawdzenie szesnastu wariantów konfiguracji hiperparametrów. Wybrane modele poddano dodatkowej walidacji na kilku ziarnach symulacji. Po zakończeniu tego etapu zamrożono model końcowy i oceniono go na tych samych scenariuszach testowych co metody referencyjne.

\subsection{Najważniejsze wyniki}
\label{subsec:podsumowanie-wyniki}

Modelem końcowym został checkpoint konfiguracji bazowej po 225\,000 kroków, wytrenowany na zbiorze o regularnym napływie pojazdów. Bardziej nieregularny zbiór treningowy nie poprawił wyników walidacyjnych, a żaden z szesnastu sprawdzonych wariantów hiperparametrów nie przewyższył konfiguracji bazowej na całym zbiorze walidacyjnym. Dalsze prowadzenie treningu do 300\,000 kroków również nie przyniosło poprawy, co potwierdziło znaczenie wyboru checkpointu na podstawie walidacji zamiast przyjmowania ostatniego zapisanego modelu.

Agent DQN uzyskał najniższe średnie opóźnienie i średni czas oczekiwania we wszystkich czterech grupach scenariuszy testowych. Średnie opóźnienie wyniosło odpowiednio 13,90~s dla ruchu lekkiego, 17,79~s dla średniego, 51,82~s dla wysokiego oraz 34,23~s dla zmiennego. Względem sterowania stałoczasowego oznacza to redukcję od 20,8\% do 50,4\%, natomiast względem sterowania akomodacyjnego od 1,3\% do 33,6\%. Największa przewaga nad metodą akomodacyjną wystąpiła przy wysokim natężeniu, a najmniejsza przy średnim, gdzie obie metody osiągnęły zbliżone rezultaty.

Przy wysokim natężeniu agent uzyskał również większą przepustowość niż obie metody referencyjne. Przewaga DQN była więc najbardziej widoczna w warunkach, w których sposób podziału czasu między fazy wpływał na zdolność obsłużenia popytu.

DQN nie był jednak najlepszy według każdej metryki. W scenariuszu zmiennym sterowanie akomodacyjne uzyskało mniejszą średnią długość kolejki, równą 8,88 pojazdu wobec 9,55 dla agenta. Metoda akomodacyjna osiągnęła także nieznacznie lepszy kwantyl rzędu 0,95 czasu oczekiwania przy lekkim i średnim ruchu. Program stałoczasowy przełączał sygnalizację znacznie rzadziej i działał bardziej przewidywalnie. W trzech z dwudziestu przebiegów modelu końcowego niewielka grupa pojazdów nie opuściła sieci przed limitem czasu, przy czym w każdym z tych przebiegów stanowiła mniej niż 1\% zgłoszonego popytu. Brak metryk tych pojazdów ogranicza siłę wniosków dotyczących przerwanych przebiegów i wskazuje aspekt wymagający dalszego doskonalenia.

\subsection{Odpowiedź na pytanie badawcze}
\label{subsec:odpowiedz-badawcza}

Na podstawie przeprowadzonego eksperymentu na pytanie badawcze można odpowiedzieć twierdząco: w przyjętych warunkach odpowiednio wytrenowany agent DQN sterował sygnalizacją skuteczniej niż badane sterowanie stałoczasowe i akomodacyjne pod względem średniego opóźnienia oraz średniego czasu oczekiwania. Przewaga wystąpiła we wszystkich grupach scenariuszy, chociaż jej skala zależała od natężenia ruchu. Przy średnim obciążeniu różnica względem sterowania akomodacyjnego była niewielka, natomiast przy wysokim natężeniu agent uzyskał wyraźnie lepsze wyniki także pod względem przepustowości.

Hipotezę badawczą uznano zatem za potwierdzoną warunkowo w zakresie badanego eksperymentu. Została ona potwierdzona dla opóźnienia i czasu oczekiwania oraz w większości scenariuszy dla długości kolejek. Wyjątek stanowił scenariusz zmienny, w którym podstawowy wynik kolejki był korzystniejszy dla sterowania akomodacyjnego. Rezultaty nie oznaczają, że DQN przewyższa każdą metodę klasyczną, w każdych warunkach i według każdego kryterium. Pokazują natomiast, że uczenie ze wzmocnieniem może doprowadzić do powstania skutecznej polityki sterowania, konkurencyjnej wobec starannie skonfigurowanych metod referencyjnych.

Zakres pracy świadomie ograniczono do pojedynczego skrzyżowania, syntetycznych scenariuszy i jednego głównego algorytmu RL. Wynik należy więc rozumieć jako potwierdzenie potencjału DQN w rozpatrywanym zadaniu, a nie jako ocenę wszystkich systemów sterowania ruchem ani gotowości rozwiązania do bezpośredniego wdrożenia. Dalsze badania powinny objąć kalibrację na danych pomiarowych, inne geometrie, więcej ziaren treningu, dodatkowych uczestników ruchu oraz pełną warstwę bezpieczeństwa.
